%-*-latex-*-
%-*-latex-*-
\newcommand\COURSE{ciss240}
\newcommand\ASSESSMENT{q1203}
\newcommand\ASSESSMENTTYPE{Quiz}
\newcommand\POINTS{\textwhite{xxx/xxx}}

\input{myquizpreamble}
\input{yliow}
\input{\COURSE}
%-*-latex-*-

%-*-latex-*-

%-*-latex-*-
\renewcommand\TITLE{\ASSESSMENTTYPE \ \ASSESSMENT}

\renewcommand\EMAIL{}
\renewcommand\AUTHOR{eagraham1@cougars.ccis.edu}

\textwidth=6in
\begin{document}
\topmatterthree

This is a closed-book, no compiler, 5 minute quiz.

%------------------------------------------------------------------------------
\nextq
Write a program that simulates the rolling of two dice until
the sum of the two dice is 10.
The program prints the dice values of each throw.
You must use a while-loop.
The output of a sample execute is given below (notice that the
last output has two integer values adding up to 10.)
Here's a test run:
\begin{console}[fontsize=\small]
1 5
2 3
5 6
6 4
\end{console}
\ANSWER
\begin{answercode}
#include <iostream>
#include <ctime>

int main()
{
    srand((unsigned int) time(NULL));
    
    while (true)
    {
        // generate a die roll and store it in die0
        int die0 = rand() % 6 + 1;
        // generate a die roll and store it in die1
        int die1 = rand() % 6 + 1;
        // print die0, die1
        std::cout << die0 << ' ' << die1 << '\n';
        // if sum of die0 and die1 is 10, break
        if(die0 + die1 == 10)
        {
             break;
        }
    }

    return 0;
}
\end{answercode}

%------------------------------------------------------------------------------
\nextq
Write a program that gets a positive integer $n$ from the user
and continually compute a sequence of integers starting with $n$.
Here's how the sequence is computed.
If a number $x$ in the sequence is even, the \textit{next} number
after $x$ is $x / 2$.
If a number $x$ in the sequence is odd, the \textit{next} number
is $3x + 1$.
Once the sequence reaches $1$, the sequence stops.
For instance suppose the integer entered by the user is $7$.
Then the sequence is $7, 22, 11, 34, 17, 52, 26, 13, 40, 20, 10, 5,
16, 8, 4, 2, 1$.
($7$ is odd, so the next number is $3 \cdot 7 = 1 = 22$.
$22$ is even, so the next number is $22/2 = 11$.
Etc.)
Here's a test run:
\begin{console}[fontsize=\footnotesize,commandchars=\\\{\}]
\userinput{7}
7 22 11 34 17 52 26 13 40 20 10 5 16 8 4 2 1
\end{console}
\ANSWER
\begin{answercode}
#include <iostream>

int main()
{
    int n;
    std::cin >> n;

    while (n != 1) 
    {
        std::cout << n << " "; 

        if (n % 2 == 0) // If the number is even
        {
            n /= 2;
        }
        else // If the number is odd
        {
            n = 3 * n + 1;
        }
    }

    std::cout << n << std::endl;
    return 0;
}
\end{answercode}
(The above procedure always produce a sequence that reaches $1$.
But the reason is not known and is a famous unsolved problem called
Collatz Conjecture.)

%------------------------------------------------------------------------------
%-*-latex-*-
\newpage
\input{instructions}
\end{document}

